\section{Existing Work and Limitations}
scMINER provides an integrated framework for clustering single-cell RNA sequencing data, reconstructing regulatory networks, and identifying transcription factor and signaling drivers \cite{pan2025scminer}. Its Mutual Information-based Clustering Analysis (MICA) is designed to capture nonlinear relationships between cells and was shown to perform well compared with commonly used approaches such as Seurat, Scanpy, SC3s, scVI, and scDeepCluster. MICA works in two modes. Below about 5,000 cells it builds the full cell by cell mutual information distance matrix, reduces it with multidimensional scaling, and runs k-means. Above that size it builds an approximate nearest neighbor graph under the same distance, embeds it with node2vec, and runs Louvain. However, several aspects of the existing work leave room for further investigation.

\subsection{The reported accuracy depends on knowing the answer in advance}
In the original benchmarking experiments, the number of clusters for MICA and the comparison methods was selected using information from the ground-truth cell labels. For larger datasets, the Louvain resolution was varied until the predicted number of clusters matched the known number of cell populations, and the result with the highest ARI was selected when multiple settings produced the same number of clusters \cite{pan2025scminer}. This provides a controlled comparison between algorithms, but it does not completely represent exploratory single-cell analysis, where the true number of cell populations is generally unknown. In the released software the user supplies the cluster count directly in the small-data mode, and the large-data mode sweeps a resolution parameter, so neither mode offers a rule for choosing a final partition without labels. This matters because the reported margin over the baselines may shrink, disappear, or move to a different method once that information is removed, and a user analyzing new tissue has no labels to tune with.

\subsection{Scalability has been measured, but accuracy at scale has not}
The original work demonstrated successful processing of datasets containing up to approximately 650,000 cells and reported competitive runtime and memory use \cite{pan2025scminer}. However, the authors also noted that calculating mutual-information-based cell similarities across genes is computationally demanding and that the absence of reliable ground-truth labels for atlas-scale datasets prevented accuracy evaluation at that scale \cite{pan2025scminer}. The largest labeled dataset in the benchmark has 13,605 cells, so the quality of mutual information clustering above that size is an open question. The software has no GPU support, and in our reading of the released code the neighbor graph is converted to a dense cell by cell matrix before clustering, which would need roughly 32 TB of memory at two million cells. Modern single-cell atlases contain millions of cells \cite{yao2023allen}, so this gap decides whether the method is usable on current data at all. The authors also report that scMINER requires at least 500 cells for some analyses, which can limit its use for rare cellular populations \cite{pan2025scminer}.

\subsection{What the mutual information between two cells actually measures is untested}
MICA places each cell's expression values into equal-width bins between that cell's own minimum and maximum, with the number of bins set by the cube root of the gene count, and estimates mutual information from the resulting joint histogram. Single-cell counts are typically 80 to 95 percent zeros, so the zero bin dominates that histogram. The score may therefore reflect mainly which genes are detected in both cells rather than how strongly they are expressed. This is not a small concern, because detection patterns alone separate cell types about as well as expression values \cite{qiu2020dropout}, and an independent comparison of similarity measures found mutual information weaker than correlation-based measures for single-cell tasks \cite{skinnider2019evaluating}. Since the bin edges come from each cell's own range, the estimate can also shift with sequencing depth. If most of the signal is co-detection, then both the reported advantage of the method and its high computational cost deserve a second look, and a better estimator becomes the natural fix.

\subsection{Cluster interpretation does not use the regulatory information the framework produces}
scMINER provides marker information, gene-network inference, and hidden-driver identification, but the framework does not evaluate a retrieval-grounded language-model approach for translating these outputs into literature-supported biological descriptions. Recent studies have shown that large language models can infer cell types from marker genes, while newer approaches use biomedical literature retrieval to improve the reliability of those annotations \cite{hou2024gpt4,crowley2025anndictionary,kazmi2025beyond,yang2026rag}. These studies primarily use marker-gene information as the biological input. Hidden drivers are exactly the evidence that marker-based annotation cannot see, since they are defined by a change in inferred activity without a matching change in expression. We have not found an annotation method that uses computed regulator activity as retrieval evidence, so whether it adds anything beyond markers is untested.
